%%%%%%%%%%%%%%%%%%%%%%%%%%%%%%%%%%%%%%%%%%%%%%%%%%%%%%%%%%%%%%%%%%%%%%%%%%%%%%%%
%
% Curriculum Vitae
% Author: Andrew Hard
% Date: November 26, 2016
%
%%%%%%%%%%%%%%%%%%%%%%%%%%%%%%%%%%%%%%%%%%%%%%%%%%%%%%%%%%%%%%%%%%%%%%%%%%%%%%%%

\documentclass{letter}

\usepackage{amssymb} % for black square
\usepackage[usenames, dvipsnames]{color}
%\definecolor{Maroon}{cmyk}{0, 0.87, 0.68, 0.32}
\definecolor{Maroon}{cmyk}{0, 0.9, 0.8, 0.55}
\usepackage{geometry} % for paper geometry
\usepackage{graphicx}
\usepackage{courier}
\usepackage{hyperref} % for hyperlinks
\hypersetup{hidelinks}
%\hypersetup{colorlinks=true,urlcolor=blue}
%\usepackage{helvetica}
\usepackage{newcent}
% Set the paper size and margin size (use a4paper, not legalpaper):
\geometry{a4paper, portrait, margin=0.5in}
\begin{document} 

%%%%%%%%%%%%%%%%%%%%%%%%%%%%%%%%%%%%%%%%%%%%%%%%%%%%%%%%%%%%%%%%%%%%%%%%%%%%%%%%
\begin{tabular}{p{0.45\textwidth}p{0.55\textwidth}}
	\hfill \newline \href{https://ch.linkedin.com/in/andrew-hard-25b690a5}{\Huge{\color{Maroon}{Andrew Hard}}} \newline \LARGE{\textit{Curriculum Vitae}} \newline
	&
	\hfill \newline 1 Hill St., Apt. 402, San Francisco, CA 94110 \newline
	(+1) 628-202-4377 \newline
	\href{mailto:hardandrew1@gmail.com}{hardandrew1@gmail.com} \newline
	\href{https://github.com/rasumovsky}{github.com/rasumovsky}\\
\end{tabular}

%%%%%%%%%%%%%%%%%%%%%%%%%%%%%%%%%%%%%%%%%%%%%%%%%%%%%%%%%%%%%%%%%%%%%%%%%%%%%%%%
\begin{flushleft}
\Large{\textsc{\textbf{\color{Maroon}{Education}}}}
\vspace{1pt} % space between text and line
\hrule
\end{flushleft}

\begin{tabular}{p{0.15\textwidth}p{0.8\textwidth}}
	2016
	&
	\textbf{Doctor of Philosophy} in Physics \newline 
	\textit{University of Wisconsin}, Madison WI, USA \newline
	Thesis: \textit{Search and discovery with the resonant $\gamma\gamma$ final state at ATLAS} \newline
	Advised by Prof. Sau Lan Wu \\
\\
	2010 
	& 
	\textbf{Bachelor of Arts} in Physics, Honors \newline 
	\textit{University of Chicago}, Chicago IL, USA \newline
	Advised by Prof. Edward Blucher 
\\
\end{tabular}

%%%%%%%%%%%%%%%%%%%%%%%%%%%%%%%%%%%%%%%%%%%%%%%%%%%%%%%%%%%%%%%%%%%%%%%%%%%%%%%%
\begin{flushleft}
\Large{\textsc{\textbf{\color{Maroon}{Experience}}}}
\vspace{1pt} % space between text and line
\hrule
\end{flushleft}

%\renewcommand{\arraystretch}{0.5}% Wider
\begin{tabular}{p{0.15\textwidth}p{0.8\textwidth}}
        2017 - Present
        &
        \textbf{Software Engineer}, \textit{Google} \newline
	\vspace{-15pt}      % suppress blank line after tabbing
        \begin{itemize}
		\itemsep0em
%		\renewcommand{\labelitemi}{\scriptsize$\blacksquare$} 
		\renewcommand{\labelitemi}{\tiny$\blacksquare$}
    		\item Developed multi-word prediction algorithms for Gboard with C++, Java
                \item Experimented with ML techniques (reinforcement learning, RNNs, LSTMs, CNNs)
        \end{itemize}
\\
	2011 - 2016
	&
	\textbf{Graduate Research Assistant}, \textit{Department of Physics, University of Wisconsin} \newline
	\vspace{-15pt}      % suppress blank line after tabbing
	\begin{itemize}
		\itemsep0em
		\renewcommand{\labelitemi}{\tiny$\blacksquare$}
    		\item Discovered Higgs boson, performed first measurements of mass, couplings, and spin 
		\item Optimized physics searches with TB-scale datasets using machine learning techniques
    		\item Developed widely used analysis software with C++, Python, ROOT, \& shell scripts
    		\item Statistical expert, created new Monte Carlo method to reduce CPU usage by $1000\times$
%		\item Invented algorithm to spatially and temporally match CMOS chip hits at LBNL 
		\item Contributed significantly to 21 papers \& notes, author on 250+ ATLAS publications
%		\item Wrote and coordinated sucessful DoE funding reports for Wisconsin ATLAS Group
	\end{itemize}
\\
	2014 
	& 
	\textbf{Graduate Teaching Assistant}, \textit{Department of Physics, University of Wisconsin} \newline
	\vspace{-15pt}      % suppress blank line after tabbing
	\begin{itemize}
		\itemsep0em
		\renewcommand{\labelitemi}{\tiny$\blacksquare$} 
		\item Led discussions and labs on classical mechanics, electrodynamics, thermodynamics
    		\item Designed supplemental exercises and summary notes that boosted exam performances
	\end{itemize}
\\
	2010 - 2011
	&
	\textbf{CERN Technologist}, \textit{Enrico Fermi Institute, University of Chicago} \newline
	\vspace{-15pt}      % suppress blank line after tabbing
	\begin{itemize}
		\itemsep0em
		\renewcommand{\labelitemi}{\tiny$\blacksquare$} 
   		\item Electronic calibration expert for the ATLAS Experiment hadronic calorimeter
   		\item Developed \& maintained calibration software package using Python and MySQL
%    		\item Documented, monitored and reported on detector status to collaboration
	\end{itemize}
\\
	2009 - 2010
	&
	\textbf{Undergraduate Research Assistant}, \textit{Enrico Fermi Institute, University of Chicago} \newline
	\vspace{-15pt}      % suppress blank line after tabbing
	\begin{itemize}
		\itemsep0em
		\renewcommand{\labelitemi}{\tiny$\blacksquare$} 
   		\item Developed particle detector simulation in C++ with ROOT and Geant4 libraries
%   		\item Constructed $\mu$ particle modules, worked in machine shop, tested electronics 
	\end{itemize}
\end{tabular}

\vspace{-10pt}


%%%%%%%%%%%%%%%%%%%%%%%%%%%%%%%%%%%%%%%%%%%%%%%%%%%%%%%%%%%%%%%%%%%%%%%%%%%%%%%%
\begin{flushleft}
\Large{\textsc{\textbf{\color{Maroon}{Skills}}}}
\vspace{1pt} % space between text and line
\hrule
\end{flushleft}

\begin{tabular}{p{0.15\textwidth}p{0.8\textwidth}}
	{\bf Scientific} 
	&
	Physics, Statistics, Simulation, Numerical Methods, Data Structures, High Throughput Computing, Databases, Machine Learning, Public Presentation \newline
\\
	{\bf Programming} 
	&
	C++, Python, Java, Go, \LaTeX{}, Unix/Linux shell scripting, ROOT, Matlab, SQL, TensorFlow \newline
\\ 
	{\bf Languages}
	&
	English (native), French (basic oral and written communication)
\end{tabular}

%%%%%%%%%%%%%%%%%%%%%%%%%%%%%%%%%%%%%%%%%%%%%%%%%%%%%%%%%%%%%%%%%%%%%%%%%%%%%%%%
\begin{flushleft}
\Large{\textsc{\textbf{\color{Maroon}{Volunteering \& Outreach}}}}
\vspace{1pt} % space between text and line
\hrule
\end{flushleft}
	\vspace{-10pt}      % suppress blank line after tabbing
\begin{tabular}{p{0.15\textwidth}p{0.8\textwidth}}
	\href{https://twiki.cern.ch/twiki/pub/AtlasPublic/HiggsPublicResults//Hgg-FloatingScale-Short2.gif}{\raisebox{-0.95\height}{\includegraphics[scale=0.16]{AnimationQR.jpg}}}
	&
	\begin{itemize}
		\itemsep0em
		\renewcommand{\labelitemi}{\tiny$\blacksquare$} 
		\item Newtonian physics demonstration for Chicago Public Library \hfill 2016
		\item US voter outreach \& registration at CERN \hfill 2016
		\item Discussed research \& funding with U.S. lawmakers in Washington D.C. \hfill 2014, 2015
		\item Created GIF visualizations of Higgs boson discovery data \hfill 2013
		\item Visited classrooms at the Chattanooga School for the Arts \& Sciences \hfill 2012
	\end{itemize}	
\end{tabular}

\vspace{-10pt}

%%%%%%%%%%%%%%%%%%%%%%%%%%%%%%%%%%%%%%%%%%%%%%%%%%%%%%%%%%%%%%%%%%%%%%%%%%%%%%%%
\begin{flushleft}
\Large{\textsc{\textbf{\color{Maroon}{Awards}}}}
\vspace{1pt} % space between text and line
\hrule
\end{flushleft}

\begin{tabular}{p{0.15\textwidth}p{0.8\textwidth}}
	2015
	&
	\textbf{Teaching Assistant Rookie of the Year}, \textit{Department of Physics, University of Wisconsin} 
\\
	2013, 2014
	& 
	\textbf{Lightning Round Winner}, \textit{US LHC User's Association Annual Meeting}
\\
\end{tabular}


%%%%%%%%%%%%%%%%%%%%%%%%%%%%%%%%%%%%%%%%%%%%%%%%%%%%%%%%%%%%%%%%%%%%%%%%%%%%%%%%
\begin{flushleft}
\Large{\textsc{\textbf{\color{Maroon}{Selected Publications}}}}
\vspace{1pt} % space between text and line
\hrule
\end{flushleft}

\textit{Search for resonances in diphoton events at $\sqrt{s}=13$ TeV with the ATLAS detector}, ATLAS Collaboration, \href{http://link.springer.com/article/10.1007/JHEP09(2016)001}{J. High Energ. Phys. (2016) 2016: 1. doi:10.1007/JHEP09(2016)001}, \href{http://arxiv.org/abs/1606.03833}{\texttt{arXiv:1606.03833 [hep-ex]}}.

\textit{Search for Higgs boson pair production in the $b \bar{b} \gamma \gamma$ final state using $pp$ collision data at $\sqrt{s}=13$ TeV with the ATLAS detector}, ATLAS Collaboration, ATLAS-CONF-2016-004, \href{https://cds.cern.ch/record/2138949}{\texttt{https://cds.cern.ch/record/2138949}}.

\textit{A search for new phenomena in events with missing $p_{T}$ and a Higgs boson decaying to two photons in a 13.3 fb$^{-1}$ $pp$ collision dataset at $\sqrt{s}=13$ TeV with the ATLAS detector}, ATLAS Collaboration, \href{https://cds.cern.ch/record/2206281}{ATLAS-CONF-2016-087}.

\textit{Performance of Silicon Pixel Detectors at Small Track Incidence Angles for the ATLAS Inner Tracker Upgrade}, ATLAS Collaboration, ATL-INDET-PROC-2015-011, \href{https://cds.cern.ch/record/2065104}{\texttt{https://cds.cern.ch/record/2065104}}.

\textit{Search for non-pointing and delayed photons in the diphoton and missing transverse momentum final state in 8 TeV $pp$ collisions at the LHC using the ATLAS detector}, ATLAS Collaboration, 
\href{http://journals.aps.org/prd/abstract/10.1103/PhysRevD.90.112005}{Phys. Rev. D90, 112005 (2014)}, \href{https://arxiv.org/abs/1409.5542}{\texttt{arXiv:1409.5542 [hep-ex]}}.

\textit{Evidence for the spin-0 nature of the Higgs boson using ATLAS data}, ATLAS Collaboration, \href{http://www.sciencedirect.com/science/article/pii/S0370269313006527?np=y}{Phys. Lett. B726 (2013) 120}, \href{https://arxiv.org/abs/1307.1432}{\texttt{arXiv:1307.1432 [hep-ex]}}.

\textit{Measurement of Higgs boson production in the diphoton decay channel in $pp$ collisions at center-of-mass energies of 7 and 8 TeV with the ATLAS detector}, ATLAS Collaboration, \href{http://journals.aps.org/prd/abstract/10.1103/PhysRevD.90.112015}{Phys. Rev. D90, 112015 (2014)}, \href{http://arxiv.org/abs/1408.7084}{\texttt{arXiv:1408.7084 [hep-ex]}}.

\textit{Observation of a new particle in the search for the Standard Model Higgs boson with the ATLAS detector at the LHC}, ATLAS Collaboration, \href{http://www.sciencedirect.com/science/article/pii/S037026931200857X}{Phys. Lett. B716 (2012) 1-29}, \href{https://arxiv.org/abs/1207.7214}{\texttt{arXiv:1207.7214 [hep-ex]}}.

\begin{flushright}
Significant contributions to 20 papers \& notes since 2011, author on \href{https://twiki.cern.ch/twiki/bin/view/AtlasPublic}{250+ ATLAS publications} since 2013.\\
%Significant contributions to 19 papers \& conference notes since 2011. \\
%Author on 250+ ATLAS publications since 2013. \\
\end{flushright}

%%%%%%%%%%%%%%%%%%%%%%%%%%%%%%%%%%%%%%%%%%%%%%%%%%%%%%%%%%%%%%%%%%%%%%%%%%%%%%%%
\begin{flushleft}
\Large{\textsc{\textbf{\color{Maroon}{Conference Presentations}}}}
\vspace{1pt} % space between text and line
\hrule
\end{flushleft}

\begin{tabular}{p{0.15\textwidth}p{0.8\textwidth}}
	August 2016
	&
	\href{https://cds.cern.ch/record/2208271}{\textbf{Search for the production of the Higgs boson in association with invisible particles in the ATLAS detector}} (Poster) \newline 
	\textit{38$^{th}$ International Conference on High Energy Physics}, Chicago, USA \\
\\
	July 2016
	&
	\href{https://cds.cern.ch/record/2199807}{\textbf{Search for a high mass diphoton resonance using the ATLAS detector}} (Invited talk) \newline 
	\textit{22$^{nd}$ International Symposium on Particles, Strings and Cosmology}, ICISE, Vietnam \newline
\\
	November 2014
	&
	\href{https://indico.hep.anl.gov/indico/contributionDisplay.py?contribId=39&sessionId=0&confId=410}{\textbf{Search for non-pointing and delayed photons in the diphoton and missing transverse momentum final state in 8 TeV $pp$ collisions}} (Talk) \newline
	\textit{2014 US LHC User's Association Annual Meeting}, Argonne National Laboratory, USA \newline
\\
	April 2014
	&
	\textbf{Higgs to diphoton workshop perspective} (Invited talk) \newline
	\textit{ATLAS Higgs Workshop}, Rome, Italy \newline
\\
	December 2013
	&
	\href{https://cds.cern.ch/record/1640386}{\textbf{Individual and combined measurements of the spin and parity properties of the Higgs boson using the ATLAS detector}} (Invited talk) \newline
	\textit{High Energy Physics in the LHC Era}, Valpara\'{i}so, Chile \newline
%\\
%	November 2013
%	&
%	\textbf{Spin analysis: method and uncertainties} (Talk) \newline
%	\textit{ATLAS Higgs to Diphoton Workshop}, Annecy, France \newline
\\
	November 2013
	&
	\textbf{Spin determination of a narrow resonance near 125 GeV with the two-photon decay channel at ATLAS} (Invited talk) \newline
	\textit{2013 US LHC User's Association Annual Meeting}, Madison, USA \newline
\\
	August 2013
	&
	\href{https://indico.bnl.gov/contributionDisplay.py?contribId=155&sessionId=9&confId=603}{\textbf{Spin measurement of the Higgs-like resonance observed in the two photon decay channel in ATLAS}} (Talk) \newline
	\textit{2013 APS Division of Particles and Fields Meeting}, SCIPP, Santa Cruz, USA \newline
%\\
%	July 2013
%	&
%	\textbf{Spin measurement of the Higgs-like resonance observed in the two-photon decay channel} (Talk) \newline
%	\textit{US ATLAS Workshop}, Argonne National Laboratory, USA \newline
\\
	November 2012
	&
	\textbf{$h\rightarrow \gamma \gamma$ vector boson fusion} (Invited talk) \newline
	\textit{US ATLAS Diboson Jamboree}, Brookhaven National Laboratory, USA \newline
%\\
%	November 2012
%	& 
%	\textbf{$VH$ lepton + $E_{T}^{miss}$} (Talk) \newline
%	\textit{ATLAS Higgs to Diphoton Workshop}, DESY, Germany
\end{tabular}
\vspace{-10pt}


%%%%%%%%%%%%%%%%%%%%%%%%%%%%%%%%%%%%%%%%%%%%%%%%%%%%%%%%%%%%%%%%%%%%%%%%%%%%%%%%
\begin{flushleft}
\Large{\textsc{\textbf{\color{Maroon}{References}}}}
\hrule
\end{flushleft}

\begin{tabular}{p{0.33\textwidth}p{0.33\textwidth}p{0.33\textwidth}}
	\textbf{Prof. Sau Lan Wu} \newline
	\textit{University of Wisconsin} \newline
	\href{mailto:Sau.Lan.Wu@cern.ch}{Sau.Lan.Wu@cern.ch} \newline
	(+41) 76 48 74 443
	&
	\textbf{Prof. John Parsons} \newline
	\textit{Columbia University} \newline
	\href{mailto:parsons@nevis.columbia.edu}{parsons@nevis.columbia.edu} \newline
	(+1) 914 591 2820
	&
	\textbf{Dr. Tancredi Carli} \newline
	\textit{CERN} \newline
	\href{mailto:Tancredi.Carli@cern.ch}{Tancredi.Carli@cern.ch} \newline
	(+41) 22 76 71 120
\end{tabular}

\end{document}